\section{External decision contract and reproducibility}
\label{sec:external-reproducibility}

\subsection{Stage-specific interpretation rules}
Table~\ref{tab:stage-rules} gives the frozen comparison and synthesis rules.
The V6 development translation rule requires all six zero-cost core sparse
cells to pass its magnitude conditions. COCO E4-D instead requires positive
gain, denominator and capture at at least one primary FORCED\_K budget.
Neither rule is inferred from the rounded display values.

\PaperTable{tables/stage_rules.tex}

\subsection{What the external classification consumes}
E1 evaluates the absolute-MAE contrast of R1 against the capacity-matched
global model; E2 evaluates the corresponding DRRE contrast. The fixed
paired intervals determine their supported, opposite or unresolved
statuses. E3 reports R1 against RECT and UNION on both metrics using
the registered practical margin. It supplies four pairwise results and
does not introduce a separate overall-support gate.

For R3, E4 requires all four frozen conditions: aggregate normalized regret
below the STOP reference; a strictly positive lower interval endpoint for
$1-\operatorname{NR}_0$; at least one predicted non-STOP decision; and
positive realized gain, positive oracle denominator and positive capture
at at least one registered sparse budget. Comparisons use unrounded saved
values. No practical-equivalence margin is inserted into these strict
utility predicates. The separate descriptive harmfulness flag is not an
E5 input.

\PaperTable{tables/coco_e4.tex}

E5 consumes the frozen E1/E2 statuses and E4-supported boolean. In the
observed branch, the fidelity predicates hold and utility support does
not, yielding \texttt{GAP\_REPLICATES}. This label states exactly that
the tested fidelity structure replicated while the defined utility
criterion did not. It does not mean that every method or budget failed,
that the true intervention effect is zero, or that a different repair
system cannot succeed. Missing or malformed inputs would produce an
execution hold, not an invented negative scientific result. E5 summarizes
these inputs; it is not another independent significance test.

\subsection{Recorded freezing and target-access sequence}
\label{sec:freeze-timeline}
Table~\ref{tab:freeze-timeline} separates prediction freezing from the later
pre-GT additions. The E4 record explicitly identifies an operationalization
amendment after predictions had been frozen; it does not claim that this
executable support rule existed at an earlier historical freeze. The
FORCED\_K primary binding was also recorded before first GT access. The E5
contract records the frozen prediction identity and pre-GT creation, but no
exact creation UTC in the bound contract or its supplement. The prediction
lock likewise supplies no exact lock UTC in the bound snapshot. We do not
replace either missing timestamp with a Git commit time.

\PaperTable{tables/freeze_timeline.tex}

Here ``registered'' refers to the repository-recorded contract at its stated
stage. It does not assert independent public preregistration. The completed
VOC FIT analysis remains development on previously used labels; the COCO
sequence establishes prediction-before-GT and pre-GT decision-contract
freezing, with the post-prediction amendments disclosed above.

\subsection{Evidence provenance and evaluation records}
The COCO records contain one ground-truth read and
\nval{coco.results} performance rows. Evaluation and scientific QA completed
before the official runner exited with code 1 because the ledger file was
absent from its incomplete working copy. The ledger was subsequently
reconstructed from its historical prefix and the unchanged append block
for the saved results, without another GT read or metric computation.
The original exit code remains 1. The two result reports match their
authenticated source snapshots; the accompanying source map records their
identities and the first-read timestamp. Earlier pre-GT records with zero
access counts describe their stage, not the completed evaluation.

Each table and figure is linked to its experiment, population, feature
schema, model, metric denominator, evidence role and source identity.
Displayed rounding leaves the underlying values unchanged, and intervals
retain their original estimands. RC-001 preserves two separately identified
sources for its main and sensitivity analyses. These records support the
reported summaries; they do not certify unavailable checkpoints or missing
source assets. The reported evidence consists of aggregate results and
derived displays; raw protected labels and access credentials are not
disclosed.
